\chapter{Capitolo 10 - Teoremi differenziali D01-D05 e applicativi A01-A03}

\section{1. Problema}

Il Capitolo 9 ha fissato il blocco fondativo F01-F10. Con questo capitolo cambiamo prospettiva: non chiediamo più soltanto ``quali leggi di base reggono la teoria?'', ma ``quanto la teoria riesce a discriminare differenze reali in casi dove il classico tende a collassare distinzioni operative?''.

Questo è il ruolo dei teoremi differenziali D01-D05:

\begin{enumerate}
\def\labelenumi{\arabic{enumi}.}
\tightlist
\item
  separare strutture formalmente simili ma ordinativamente divergenti;
\item
  classificare perdite informative e simmetrie non universali;
\item
  mostrare il confine tra dualità formale e validità reale.
\end{enumerate}

A valle, i teoremi applicativi A01-A03 traducono la differenza teorica in protocolli valutabili su pipeline IA, ricostruzione strutturale e sistemi narrativi complessi.

Il problema centrale del capitolo è quindi:

\textbf{come trasformare la non-ridondanza teorica di OCT in capacità discriminante misurabile e replicabile?}

La difficoltà è pratica e metodologica:

\begin{enumerate}
\def\labelenumi{\arabic{enumi}.}
\tightlist
\item
  un risultato differenziale può essere formalmente elegante ma empiricamente fragile;
\item
  un risultato applicativo può essere empiricamente suggestivo ma teoricamente ambiguo.
\end{enumerate}

Per questo il capitolo adotta una regola unitaria: ogni teorema D/A deve mantenere un ponte esplicito tra forma, metrica e criterio di fallimento.

\begin{center}\rule{0.5\linewidth}{0.5pt}\end{center}

\section{2. Tesi del capitolo}

La tesi è articolata in sette punti.

\begin{enumerate}
\def\labelenumi{\arabic{enumi}.}
\tightlist
\item
  I teoremi D01-D05 mostrano che OCT distingue classi di fenomeni che il quadro classico tende a identificare.
\item
  La differenza ordinativa emerge soprattutto in presenza di sterilità emergenziale, perdita informativa e dipendenza contestuale.
\item
  Le simmetrie formali rilevanti nel classico diventano condizionate in OCT.
\item
  I teoremi applicativi A01-A03 sono estensioni operative dei teoremi fondativi/differenziali e non moduli indipendenti.
\item
  Ogni teorema applicativo richiede protocollo, metrica e criterio di confutazione esplicito.
\item
  Il capitolo non chiude definitivamente tutte le prove, ma stabilisce una forma canonicamente validabile.
\item
  Il blocco D/A è sufficiente per aprire il Capitolo 11 su controesempi, limiti e falsificabilità globale.
\end{enumerate}

\begin{center}\rule{0.5\linewidth}{0.5pt}\end{center}

\section{3. Definizioni canoniche}

\subsection{Definizione 10.1 (Differenza strutturale forte)}

Due strutture \texttt{X} e \texttt{Y} hanno differenza strutturale forte in \texttt{Omega} se:

\begin{enumerate}
\def\labelenumi{\arabic{enumi}.}
\tightlist
\item
  risultano equivalenti o indistinguibili secondo criteri classici selezionati;
\item
  divergono su almeno un invariante ordinativo chiave oltre soglia.
\end{enumerate}

\subsection{Definizione 10.2 (Perdita ontologica)}

Data una trasformazione dimenticante \texttt{U}, la perdita è:

\begin{enumerate}
\def\labelenumi{\arabic{enumi}.}
\tightlist
\item
  \textbf{preservativa} se non altera classificazione A/B/C in \texttt{S\_Omega};
\item
  \textbf{degenerativa} se induce passaggio verso sterilità o degenerazione.
\end{enumerate}

\subsection{Definizione 10.3 (Simmetria condizionata)}

Una simmetria è condizionata quando vale solo su sottofamiglie di contesti \texttt{Omega} e non in modo universale.

\subsection{Definizione 10.4 (Dualità condizionata)}

La dualizzazione formale è condizionata ordinativamente quando non preserva automaticamente \texttt{OCT\_Valid\_Omega}.

\subsection{Definizione 10.5 (Protocollo applicativo minimo)}

Un teorema applicativo è in stato forte solo se include:

\begin{enumerate}
\def\labelenumi{\arabic{enumi}.}
\tightlist
\item
  dataset o benchmark dichiarati;
\item
  metrica ordinativa esplicita (\texttt{Coh}, \texttt{Phi}, \texttt{Delta});
\item
  baseline comparativa;
\item
  criterio di accettazione e criterio di fallimento.
\end{enumerate}

\begin{center}\rule{0.5\linewidth}{0.5pt}\end{center}

\section{4. Sviluppo formale}

\subsection{4.1 D01 - Teorema di differenza strutturale forte}

\textbf{Enunciato.}\\
Esistono strutture indistinguibili nel quadro classico ma distinguibili in OCT.

Ipotesi:

\begin{enumerate}
\def\labelenumi{\arabic{enumi}.}
\tightlist
\item
  coppia \texttt{X,Y} con equivalenza classica sul profilo selezionato;
\item
  disponibilità di invarianti ordinativi comparabili;
\item
  soglia di divergenza dichiarata.
\end{enumerate}

Tesi:
\texttt{Eq\_class(X,Y)} non implica \texttt{Eq\_ord,Omega(X,Y)}.

Proof sketch:

\begin{enumerate}
\def\labelenumi{\arabic{enumi}.}
\tightlist
\item
  costruzione di una coppia classically-equivalent;
\item
  calcolo invarianti \texttt{Coh/Phi/Delta};
\item
  verifica divergenza oltre tolleranza su almeno un asse.
\end{enumerate}

Conseguenze:

\begin{enumerate}
\def\labelenumi{\arabic{enumi}.}
\tightlist
\item
  concretizza F01/F07 in forma differenziale;
\item
  fornisce template per test di non-ridondanza su domini diversi.
\end{enumerate}

Stato: \texttt{in\_proof}.

\subsection{4.2 D02 - Teorema di commutatività non produttiva}

\textbf{Enunciato.}\\
Esistono diagrammi commutativi classici con \texttt{Phi=0}.

Ipotesi:

\begin{enumerate}
\def\labelenumi{\arabic{enumi}.}
\tightlist
\item
  diagramma commutativo \texttt{D};
\item
  composizione formalmente corretta;
\item
  misura emergenziale su output compositivo.
\end{enumerate}

Tesi:
commutatività classica non implica produttività ordinativa.

Proof sketch:

\begin{enumerate}
\def\labelenumi{\arabic{enumi}.}
\tightlist
\item
  selezione di diagramma formalmente commutativo;
\item
  verifica \texttt{Phi\_Omega(D)=0};
\item
  classificazione in regione B (sterilità).
\end{enumerate}

Conseguenze:

\begin{enumerate}
\def\labelenumi{\arabic{enumi}.}
\tightlist
\item
  smonta l'inferenza ``commutativo quindi fecondo'';
\item
  rafforza il gate anti-sterilità F08.
\end{enumerate}

Stato: \texttt{in\_proof}.

\subsection{4.3 D03 - Teorema di perdita ontologica dei funtori dimenticanti}

\textbf{Enunciato.}\\
La perdita informativa dei funtori dimenticanti è classificabile in preservativa versus degenerativa.

Ipotesi:

\begin{enumerate}
\def\labelenumi{\arabic{enumi}.}
\tightlist
\item
  funtore dimenticante \texttt{U} tra strutture con layer ordinativo;
\item
  criterio di confronto pre/post trasformazione su \texttt{S\_Omega}.
\end{enumerate}

Tesi:
la perdita non è monolitica: alcune riduzioni preservano validità, altre la degradano.

Proof sketch:

\begin{enumerate}
\def\labelenumi{\arabic{enumi}.}
\tightlist
\item
  definizione di indici di perdita su \texttt{Coh} e \texttt{Phi};
\item
  partizione dei casi in preservativi e degenerativi;
\item
  verifica su esempi differenziati.
\end{enumerate}

Conseguenze:

\begin{enumerate}
\def\labelenumi{\arabic{enumi}.}
\tightlist
\item
  fornisce tassonomia tecnica della perdita ontologica;
\item
  guida progettazione di pipeline che usano proiezioni o compressioni.
\end{enumerate}

Stato: \texttt{in\_proof}.

\subsection{4.4 D04 - Teorema di simmetria condizionata}

\textbf{Enunciato.}\\
La simmetria monoidale ordinativa è dominio-dipendente e non universalmente imponibile.

Ipotesi:

\begin{enumerate}
\def\labelenumi{\arabic{enumi}.}
\tightlist
\item
  classe di domini con vincoli asimmetrici;
\item
  struttura monoidale formalmente definita;
\item
  valutazione ordinativa pre/post simmetrizzazione.
\end{enumerate}

Tesi:
la simmetria formale può rompere validità ordinativa in domini con direzionalità causale o funzionale.

Proof sketch:

\begin{enumerate}
\def\labelenumi{\arabic{enumi}.}
\tightlist
\item
  esibizione di famiglia di domini asimmetrici;
\item
  confronto tra versione simmetrica e non simmetrica;
\item
  verifica di perdita su \texttt{Coh} o \texttt{Phi} nella versione simmetrica.
\end{enumerate}

Conseguenze:

\begin{enumerate}
\def\labelenumi{\arabic{enumi}.}
\tightlist
\item
  limita l'estensione indiscriminata di intuizioni monoidali classiche;
\item
  introduce un criterio di ammissibilità della simmetria.
\end{enumerate}

Stato: \texttt{in\_proof}.

\subsection{4.5 D05 - Teorema di dualità condizionata}

\textbf{Enunciato.}\\
La dualizzazione formale non preserva sempre validità ordinativa.

Ipotesi:

\begin{enumerate}
\def\labelenumi{\arabic{enumi}.}
\tightlist
\item
  struttura \texttt{C} e duale \texttt{C\^{}op};
\item
  metrica ordinativa coerente su entrambe;
\item
  criterio di confronto su invarianti.
\end{enumerate}

Tesi:
\texttt{Dual\_formale(C)} non implica \texttt{Dual\_ord,Omega(C)}.

Proof sketch:

\begin{enumerate}
\def\labelenumi{\arabic{enumi}.}
\tightlist
\item
  costruzione di caso in cui inversione dei morfismi resta lecita sintatticamente;
\item
  osservazione di divergenza ordinativa su composizioni dualizzate;
\item
  classificazione della dualità come condizionata.
\end{enumerate}

Conseguenze:

\begin{enumerate}
\def\labelenumi{\arabic{enumi}.}
\tightlist
\item
  evita inferenze automatiche da proprietà duali classiche;
\item
  prepara la discussione dei limiti metateorici nel Capitolo 11.
\end{enumerate}

Stato: \texttt{revise}.

\subsection{4.6 A01 - Teorema di stabilità ordinativa in pipeline IA}

\textbf{Enunciato.}\\
Pipeline compositive con coerenza ordinativa alta mostrano minore degenerazione semantica rispetto a pipeline solo formalmente corrette.

Dipendenze principali:
F03, D03.

Protocollo minimo:

\begin{enumerate}
\def\labelenumi{\arabic{enumi}.}
\tightlist
\item
  benchmark multi-step di trasformazione semantica;
\item
  confronto tra pipeline baseline classica e pipeline con gate ordinativo;
\item
  misure su \texttt{Coh}, \texttt{Phi}, drift semantico, errore cumulativo.
\end{enumerate}

Tesi operativa:
con \texttt{Coh} mantenuta sopra soglia e controllo perdita ontologica, il drift cresce più lentamente.

Criterio di fallimento:
se la pipeline ordinativa non mostra vantaggio rispetto a baseline su almeno due metriche principali.

Stato: \texttt{in\_proof}.

\subsection{4.7 A02 - Teorema di ricostruzione strutturale da proiezioni linguistiche}

\textbf{Enunciato.}\\
In condizioni di aggiunzione controllata, è possibile recuperare struttura ordinativa da output linguistico meglio del baseline classico.

Dipendenze principali:
F05, F06.

Protocollo minimo:

\begin{enumerate}
\def\labelenumi{\arabic{enumi}.}
\tightlist
\item
  task di inversione testo-\textgreater struttura su dataset annotato;
\item
  stima fedeltà strutturale e consistenza inter-osservatore;
\item
  confronto con ricostruzione classica non ordinativa.
\end{enumerate}

Tesi operativa:
l'uso del gate ordinativo in fase di ricostruzione riduce soluzioni formalmente plausibili ma funzionalmente sterili.

Criterio di fallimento:
assenza di miglioramento statisticamente robusto su fedeltà e consistenza.

Stato: \texttt{in\_proof}.

\subsection{4.8 A03 - Teorema di degenerazione nei sistemi sociali narrativi}

\textbf{Enunciato.}\\
Sistemi narrativamente coerenti ma con \texttt{Phi} nulla mostrano dinamiche di controllo/degrado prevedibili.

Dipendenze principali:
F08, D02.

Protocollo minimo:

\begin{enumerate}
\def\labelenumi{\arabic{enumi}.}
\tightlist
\item
  analisi longitudinale di reti discorsive;
\item
  misura della divergenza tra coerenza retorica e funzione emergente;
\item
  tracciamento di indicatori di irrigidimento/chiusura sistemica.
\end{enumerate}

Tesi operativa:
alta coerenza narrativa senza emergenza funzionale è segnale precoce di degenerazione.

Criterio di fallimento:
assenza di relazione sistematica tra sterilità emergenziale e pattern degenerativi osservati.

Stato: \texttt{in\_proof}.

\subsection{4.9 Coerenza del blocco D/A}

Il blocco differenziale/applicativo è coerente se vale la seguente catena:

\begin{enumerate}
\def\labelenumi{\arabic{enumi}.}
\tightlist
\item
  F01-F10 stabiliscono il quadro;
\item
  D01-D05 stressano il quadro su differenze strutturali;
\item
  A01-A03 trasferiscono la differenza in protocolli empirici.
\end{enumerate}

Questa progressione impedisce due derive:

\begin{enumerate}
\def\labelenumi{\arabic{enumi}.}
\tightlist
\item
  applicazioni senza fondazione;
\item
  fondazione senza conseguenze verificabili.
\end{enumerate}

\subsection{4.10 Matrice di maturità}

Stato attuale sintetico:

\begin{enumerate}
\def\labelenumi{\arabic{enumi}.}
\tightlist
\item
  \texttt{in\_proof}: D01, D02, D03, D04, A01, A02, A03;
\item
  \texttt{revise}: D05;
\item
  \texttt{validated}: nessuno nel blocco D/A in forma completa.
\end{enumerate}

Interpretazione:

\begin{enumerate}
\def\labelenumi{\arabic{enumi}.}
\tightlist
\item
  il blocco è avanzato ma ancora in fase di consolidamento;
\item
  il valore del capitolo è fissare forma canonica, non dichiarare chiusure premature.
\end{enumerate}

\subsection{4.11 Criteri di confutazione per il blocco D/A}

Per evitare che i teoremi differenziali e applicativi restino solo descrittivi, fissiamo criteri di confutazione sintetici.

\begin{enumerate}
\def\labelenumi{\arabic{enumi}.}
\tightlist
\item
  \textbf{D01 è confutato} se ogni coppia classically-equivalent risulta anche ordinativamente equivalente senza eccezioni rilevanti.
\item
  \textbf{D02 è confutato} se non si trova alcun diagramma commutativo con \texttt{Phi=0}.
\item
  \textbf{D03 è confutato} se tutte le perdite da funtori dimenticanti risultano sempre preservative.
\item
  \textbf{D04 è confutato} se la simmetria monoidale risulta ordinativamente valida in ogni dominio testato senza condizioni.
\item
  \textbf{D05 è confutato} se la dualizzazione preserva sempre validità ordinativa in ogni classe rilevante.
\item
  \textbf{A01 è confutato} se la pipeline ordinativa non migliora rispetto al baseline su metriche chiave.
\item
  \textbf{A02 è confutato} se la ricostruzione ordinativa non supera il baseline in fedeltà e consistenza.
\item
  \textbf{A03 è confutato} se non emerge relazione tra sterilità emergenziale e dinamiche degenerative.
\end{enumerate}

Questa sezione non pretende di esaurire tutti i casi di confutazione, ma fornisce una base comune per progettare benchmark orientati a falsificazione e non solo a conferma.

\begin{center}\rule{0.5\linewidth}{0.5pt}\end{center}

\section{5. Proposizioni e teoremi locali}

\subsection{Proposizione 10.1 (Dipendenza non arbitraria dei teoremi applicativi)}

Enunciato:
Ogni teorema applicativo Axx dipende da almeno un teorema fondativo o differenziale esplicito.

Intuizione:
A01 dipende da F03/D03, A02 da F05/F06, A03 da F08/D02.

Stato: \texttt{validated}.

\subsection{Proposizione 10.2 (Sufficienza della metrica triadica per il blocco D/A)}

Enunciato:
La triade (\texttt{Coh},\texttt{Phi},\texttt{Delta}) è sufficiente come base minima di misura per i teoremi D/A.

Intuizione:
consente di distinguere stabilità, sterilità e degenerazione senza imporre metriche dominio-specifiche uniche.

Stato: \texttt{in\_review}.

\subsection{Teorema 10.3 (Trasferibilità operativa condizionata)}

Ipotesi:

\begin{enumerate}
\def\labelenumi{\arabic{enumi}.}
\tightlist
\item
  due domini con mapping contestuale esplicito;
\item
  protocolli metrici compatibili;
\item
  criteri di soglia dichiarati.
\end{enumerate}

Tesi:
I risultati D/A sono trasferibili in forma condizionata, non automaticamente universale.

Proof sketch:

\begin{enumerate}
\def\labelenumi{\arabic{enumi}.}
\tightlist
\item
  la trasferibilità richiede coerenza semantica delle variabili osservate;
\item
  con mapping esplicito, la classificazione A/B/C/D resta confrontabile;
\item
  senza mapping, i risultati non sono direttamente trasferibili.
\end{enumerate}

Conseguenze:
fornisce regola di prudenza per generalizzazioni inter-dominio.

Stato: \texttt{in\_proof}.

\begin{center}\rule{0.5\linewidth}{0.5pt}\end{center}

\section{6. Implicazioni operative}

\subsection{6.1 Per il Capitolo 11}

Il Capitolo 11 dovrà usare D01-D05 e A01-A03 come base per:

\begin{enumerate}
\def\labelenumi{\arabic{enumi}.}
\tightlist
\item
  costruire controesempi realistici, non solo astratti;
\item
  esplicitare limiti teorici e limiti empirici separatamente;
\item
  definire confini di falsificabilità robusti.
\end{enumerate}

\subsection{6.2 Per la pipeline di validazione}

Ogni teorema D/A può essere associato a un pacchetto di test dedicato:

\begin{enumerate}
\def\labelenumi{\arabic{enumi}.}
\tightlist
\item
  \texttt{Dxx\_tests} per controesempi strutturali;
\item
  \texttt{Axx\_tests} per benchmark applicativi;
\item
  report con stato e motivazione di transizione.
\end{enumerate}

Questo approccio rende il programma cumulativo e auditabile.

\subsection{6.3 Per pubblicazione e release note}

Nelle release pubbliche conviene includere una tabella standard:

\begin{enumerate}
\def\labelenumi{\arabic{enumi}.}
\tightlist
\item
  teorema;
\item
  stato;
\item
  ultima evidenza aggiunta;
\item
  gap residuo per \texttt{validated}.
\end{enumerate}

In questo modo, il lettore non confonde maturità editoriale del capitolo con maturità completa dei singoli teoremi.

\subsection{6.4 Per uso interdisciplinare}

Il blocco D/A facilità il dialogo tra matematica, IA e scienze sociali perché fornisce:

\begin{enumerate}
\def\labelenumi{\arabic{enumi}.}
\tightlist
\item
  formalismo comune;
\item
  protocolli osservabili;
\item
  linguaggio dei limiti esplicito.
\end{enumerate}

La condizione è mantenere la distinzione tra:

\begin{enumerate}
\def\labelenumi{\arabic{enumi}.}
\tightlist
\item
  prova formale;
\item
  evidenza empirica;
\item
  interpretazione di dominio.
\end{enumerate}

\subsection{6.5 Piano immediato di consolidamento}

Ordine consigliato di consolidamento:

\begin{enumerate}
\def\labelenumi{\arabic{enumi}.}
\tightlist
\item
  D02 e A01 (priorità alta per impatto dimostrativo);
\item
  D03 e A03 (perdita ontologica e degenerazione narrativa);
\item
  D01, D04, D05, A02 (raffinamenti strutturali e trasferibilità).
\end{enumerate}

Questa sequenza massimizza il rendimento scientifico nel breve periodo e prepara bene il passaggio a \texttt{v1.2} completa.

\subsection{\texorpdfstring{6.6 Criteri minimi per passaggio D/A a \texttt{validated}}{6.6 Criteri minimi per passaggio D/A a validated}}

Per coerenza con il capitolo precedente, proponiamo cinque requisiti minimi anche per D01-D05 e A01-A03.

\begin{enumerate}
\def\labelenumi{\arabic{enumi}.}
\tightlist
\item
  enunciato formalmente non ambiguo, con ipotesi dichiarate;
\item
  proof sketch riproducibile o protocollo applicativo completo;
\item
  almeno un controesempio limite o test di fallimento esplicito;
\item
  confronto con baseline classica documentato;
\item
  aggiornamento claim ledger con motivazione del cambio di stato.
\end{enumerate}

Nei teoremi applicativi aggiungiamo un requisito ulteriore: disponibilità dei dati o del manifesto benchmark necessario alla replica indipendente.

Questo vincolo è fondamentale per mantenere allineamento tra qualità teorica del manoscritto e qualità scientifica della sua disseminazione pubblica.

In assenza di questi criteri, lo stato \texttt{validated} perde significato cumulativo e diventa una semplice etichetta editoriale. L'obiettivo, invece, è costruire una progressione verificabile nel tempo, compatibile con audit esterni e con eventuali revisioni critiche della teoria.

\begin{center}\rule{0.5\linewidth}{0.5pt}\end{center}

\section{7. Limiti e non-claim}

Limiti espliciti:

\begin{enumerate}
\def\labelenumi{\arabic{enumi}.}
\tightlist
\item
  il capitolo non contiene prove complete chiuse per l'intero blocco D/A;
\item
  le metriche applicative richiedono ancora taratura per domini specifici;
\item
  D05 necessità ulteriore formalizzazione su categorie opposte in contesti non simmetrici;
\item
  A01-A03 sono in stato pre-consolidamento empirico.
\end{enumerate}

Failure mode riconosciuti:

\begin{enumerate}
\def\labelenumi{\arabic{enumi}.}
\tightlist
\item
  confondere correlazione empirica con dimostrazione teorematica;
\item
  usare benchmark troppo omogenei e sovrastimare trasferibilità;
\item
  dichiarare generalizzazioni cross-domain senza mapping contestuale;
\item
  ignorare i criteri di fallimento e riportare solo risultati positivi.
\end{enumerate}

Box C - Non-claim

Questo capitolo non implica:

\begin{enumerate}
\def\labelenumi{\arabic{enumi}.}
\tightlist
\item
  che il blocco D/A sia già completamente validato;
\item
  che ogni applicazione OCT mostri automaticamente vantaggi rispetto al classico;
\item
  che i protocolli qui descritti sostituiscano il lavoro di replica indipendente.
\end{enumerate}

\begin{center}\rule{0.5\linewidth}{0.5pt}\end{center}

\section{8. Claim Ledger (S0-S3)}

{\def\LTcaptype{none} % do not increment counter
\begin{longtable}[]{@{}lllll@{}}
\toprule\noalign{}
ID & Claim & Tipo & Evidenza & Stato \\
\midrule\noalign{}
\endhead
\bottomrule\noalign{}
\endlastfoot
C10.1 & D01-D05 estendono in modo non ridondante il blocco fondativo su scenari differenziali & formale & S2 & in\_review \\
C10.2 & D02 mostra la possibilità di commutatività classica non produttiva & formale & S2 & in\_proof \\
C10.3 & D03 consente una tassonomia operativa della perdita ontologica & formale-operativo & S2 & in\_proof \\
C10.4 & D04 e D05 limitano l'uso indiscriminato di simmetria e dualità formali & metateorico & S2 & revise \\
C10.5 & A01-A03 traducono il differenziale teorico in protocolli empirico-formali replicabili & metodologico & S1 & in\_review \\
C10.6 & La triade (\texttt{Coh},\texttt{Phi},\texttt{Delta}) è base minima sufficiente per il monitoraggio del blocco D/A & metodologico & S1 & in\_review \\
C10.7 & La trasferibilità inter-dominio è condizionata e richiede mapping contestuale esplicito & epistemico & S1 & validated \\
\end{longtable}
}

\begin{center}\rule{0.5\linewidth}{0.5pt}\end{center}

\section{9. Bridge al Capitolo 11}

Il Capitolo 10 ha trasformato il blocco fondativo in differenziale operativo:

\begin{enumerate}
\def\labelenumi{\arabic{enumi}.}
\tightlist
\item
  ha mostrato dove il classico e OCT divergono in modo misurabile;
\item
  ha fissato protocolli applicativi con criteri di fallimento;
\item
  ha chiarito lo stato di maturità reale del programma D/A.
\end{enumerate}

Il Capitolo 11 userà questo materiale per chiudere il triangolo critico della teoria:

\begin{enumerate}
\def\labelenumi{\arabic{enumi}.}
\tightlist
\item
  controesempi sistematici;
\item
  limiti strutturali ed empirici;
\item
  criteri di falsificabilità forti e condizioni di revisione.
\end{enumerate}

In sintesi: il Capitolo 10 misura la capacità discriminante di OCT; il Capitolo 11 ne prova la tenuta critica.

\begin{center}\rule{0.5\linewidth}{0.5pt}\end{center}

\section{Riferimenti}

Riferimenti di framework interno:

\begin{enumerate}
\def\labelenumi{\arabic{enumi}.}
\tightlist
\item
  Ghioni, F. (2026). \texttt{TE\_CORE\_v5.1.md}.
\item
  Ghioni, F. (2026). \texttt{Ordinative\_Set\_Theory\_OST\_A\_Concise\_Guide\_For\_AI\_v2\_1.md}.
\item
  \texttt{OCT\_THEOREM\_PROGRAM\_v0\_1.md}.
\item
  \texttt{OCT\_TYPED\_FORMAL\_SPEC\_v0\_1.md}.
\item
  \texttt{OCT\_FOUNDATIONAL\_BOOK\_CHAPTER\_09\_v1\_0.md}.
\end{enumerate}

Riferimenti primari:

\begin{enumerate}
\def\labelenumi{\arabic{enumi}.}
\tightlist
\item
  Mac Lane, S. (1998). \emph{Categories for the Working Mathematician}.
\item
  Riehl, E. (2016). \emph{Category Theory in Context}.
\item
  Spivak, D. (2014). \emph{Category Theory for the Sciences}.
\end{enumerate}
